\documentclass[10pt]{beamer}

\usetheme[numbering=none]{metropolis}
\usepackage{appendixnumberbeamer}

\usepackage[utf8]{inputenc}

\usepackage{booktabs}
\usepackage[scale=2]{ccicons}

\usepackage{verbatim}

\usepackage{pgfplots}
\usepgfplotslibrary{dateplot}

\usepackage{xspace}
\newcommand{\themename}{\textbf{\textsc{metropolis}}\xspace}

\usepackage{xcolor}

\usepackage{listings}

\usepackage{pbox}

\usepackage[absolute,overlay]{textpos}



\title{Valaris}
\subtitle{Echtzeit Computergrafik in der Spieleentwicklung}
\date{26. Juni 2018}
\author{Sidney Hansen}
\titlegraphic{\hfill\includegraphics[height=2.5cm]{./Bilder/ValarisLogo.png}}

\begin{document}

\maketitle

\begin{frame}[fragile]{Module}

    \begin{itemize}
        \item \textcolor{violet}{menu}
        \item \textcolor{green}{generation}
        \item Rendering
        \begin{itemize}
            \item \textcolor{orange}{tick}
            \item \textcolor{orange}{rendering.tick}
            \item \textcolor{red}{rendering.effects}
            \item \textcolor{red}{rendering.animation}
            \item \textcolor{red}{rendering.postprocessing}
        \end{itemize}
        \item \textcolor{blue}{assets}
    \end{itemize}

\end{frame}







\begin{frame}[fragile]{Menu}
 \setlength{\TPHorizModule}{\textwidth}
 \setlength{\TPVertModule}{\textwidth}
 \setbeamercolor{block title}{use=structure,fg=black,bg=white!75!gray}
 \setbeamercolor{block body}{use=structure,fg=black,bg=white!90!gray}

        \begin{textblock}{0.5}(0.35,0.1)
                \begin{block}{Steuerung}
                    \textit{menu.gui}
                    \begin{itemize}
                        \item KeyInputHandler
                        \item ScreenController
                        \item NiftyAppState
                    \end{itemize}
                \end{block}
        \end{textblock}
\end{frame}

\begin{frame}[fragile]{Menu}
    \setlength{\TPHorizModule}{\textwidth}
    \setlength{\TPVertModule}{\textwidth}
    \setbeamercolor{block title}{use=structure,fg=black,bg=white!75!gray}
    \setbeamercolor{block body}{use=structure,fg=black,bg=white!90!gray}
   
    \begin{textblock}{0.5}(0.35,0.1)
            \begin{block}{Steuerung}
                \textit{menu.gui}
                \begin{itemize}
                    \item KeyInputHandler
                    \item ScreenController
                    \item NiftyAppState
                \end{itemize}
            \end{block}
    \end{textblock}

    \begin{textblock}{0.5}(0.05,0.5)
        \begin{block}{Anzeige}
            \textit{de.lessvoid.nifty}
            \begin{itemize}
                \item Layout XML
            \end{itemize}
        \end{block}
    \end{textblock}

\end{frame}

\begin{frame}[fragile]{Menu}
    \setlength{\TPHorizModule}{\textwidth}
    \setlength{\TPVertModule}{\textwidth}
    \setbeamercolor{block title}{use=structure,fg=black,bg=white!75!gray}
    \setbeamercolor{block body}{use=structure,fg=black,bg=white!90!gray}
   
    \begin{textblock}{0.5}(0.35,0.1)
            \begin{block}{Steuerung}
                \textit{menu.gui}
                \begin{itemize}
                    \item KeyInputHandler
                    \item ScreenController
                    \item NiftyAppState
                \end{itemize}
            \end{block}
    \end{textblock}

    \begin{textblock}{0.5}(0.65,0.5)
            \begin{block}{Spieldaten}
                \textit{menu.menuconfig}
                \begin{itemize}
                    \item InternalGameConfig
                    \item SeedConfig
                    \item GraphicsConfig
                \end{itemize}
            \end{block}
    \end{textblock}

    \begin{textblock}{0.5}(0.05,0.5)
            \begin{block}{Anzeige}
                \textit{de.lessvoid.nifty}
                \begin{itemize}
                    \item Layout XML
                \end{itemize}
            \end{block}
    \end{textblock}
\end{frame}



\begin{frame}[fragile]{Menu}
    \setlength{\TPHorizModule}{\textwidth}
    \setlength{\TPVertModule}{\textwidth}
    \setbeamercolor{block title}{use=structure,fg=black,bg=white!75!gray}
    \setbeamercolor{block body}{use=structure,fg=black,bg=white!90!gray}
   
    \begin{textblock}{0.5}(0.35,0.1)
            \begin{block}{\textcolor{red}{Controller}}
                \textit{menu.gui}
                \begin{itemize}
                    \item KeyInputHandler
                    \item ScreenController
                    \item NiftyAppState
                \end{itemize}
            \end{block}
    \end{textblock}

    \begin{textblock}{0.5}(0.65,0.5)
            \begin{block}{\textcolor{red}{Model}}
                \textit{menu.menuconfig}
                \begin{itemize}
                    \item InternalGameConfig
                    \item SeedConfig
                    \item GraphicsConfig
                \end{itemize}
            \end{block}
    \end{textblock}

    \begin{textblock}{0.5}(0.05,0.5)
            \begin{block}{\textcolor{red}{View}}
                \textit{de.lessvoid.nifty}
                \begin{itemize}
                    \item Layout XML
                \end{itemize}
            \end{block}
    \end{textblock}
\end{frame}



\begin{frame}[fragile]{Generierung}
    \begin{itemize}
        \item[]
        \item[]
        \item[]
        \item[]
        \item[]
        \item[]
        \item \textcolor{red}{\textbf{roadgeneration:}\\Generierung der Strecke in den Segmenten}
    \end{itemize}
\end{frame}

\begin{frame}[fragile]{Generierung}
    \includegraphics[width=.3\linewidth]{./Bilder/domegenerationPrototype.PNG}
    \includegraphics[width=.3\linewidth]{./Bilder/tunnelgenerationPrototype.PNG}
    \includegraphics[width=.3\linewidth]{./Bilder/roomgenerationPrototype.PNG}
    \begin{itemize}
        \item \textcolor{red}{\textbf{domegeneration, tunnelgeneration, roomgeneration:}\\Generierung von Streckensegmenten}\par
        \item \textcolor{gray}{\textbf{roadgeneration:}\\Generierung der Strecke in den Segmenten}
    \end{itemize}
\end{frame}

\begin{frame}[fragile]{Generierung}
    
    \includegraphics[width=.5\linewidth]{./Bilder/mapgenerationPrototype.PNG}
    \begin{itemize}
        \item \textcolor{red}{\textbf{mapgeneration:}\\Zusammensetzung der Streckensegmenten}\par
        \textcolor{gray}{\item \textbf{domegeneration, tunnelgeneration, roomgeneration:}\\Generierung von Streckensegmenten}\par
        \textcolor{gray}{\item \textbf{mapgeneration:}\\Generierung der Strecke in den Segmenten}
    \end{itemize}
\end{frame}



\begin{frame}[fragile]{Generierung}
    \begin{lstlisting}
generate(...) {
    ...
    DomeGenerator domeGenerator = new DomeGenerator(...);
    Map map = new Map(...);
    ...
    Dome dome = domeGenerator.generateDome(...);
    map.getRoad.attach(dome.getRoad);
    map.addSceneItem(dome);
    ...
    rootNode.attachChild(map.generateSceneGraph());
    return map.getMapBody();
}
    \end{lstlisting}
\end{frame}




\begin{frame}[fragile]{Rendering}
    \begin{itemize}
        \item Manipulation von SceneGraph
        \begin{itemize}
            \item \textcolor{orange}{tick}
            \item \textcolor{orange}{rendering.tick}
            \item \textcolor{red}{rendering.effects}
            \item \textcolor{red}{rendering.animation}
        \end{itemize}
        \item Render Pipeline
        \begin{itemize}
            \item \textcolor{red}{rendering.postprocessing}
            \item materialdefinitionen
            \item shader
        \end{itemize}
    \end{itemize}
\end{frame}


\begin{frame}[fragile]{Rendering: Tick und Tickverarbeitung}
    Ticker
    \begin{itemize}
        \item \textbf{simulationTick:} Aktuelle Simulation
        \item \textbf{currentTick:} Zwischenspeicher
        \item \textbf{renderBuffer:} Wird aktuell gerendert
    \end{itemize}
\end{frame}

\begin{frame}[fragile]{Rendering: Tick und Tickverarbeitung}
    Simulation
    \begin{itemize}
        \item \textbf{Tick:} Gesammtzustand aller Dynamischen Objekte
        \begin{itemize}
            \item \textit{TickEvent}
        \end{itemize}
        \item \textbf{DynamicGameObject:} Zustand eines Dynamischen Objektes
        \begin{itemize}
            \item \textit{DynamicGameObjectEvent}
        \end{itemize}
        \item \textbf{IProperty:} Eine Eigenschaft eines Dynamischen Objektes
    \end{itemize}

    \mbox{}\par
    Render
    \begin{itemize}
        \item \textbf{TickProcessor}
        \begin{itemize}
            \item \textit{TickEventProcessor}
        \end{itemize}
        \item \textbf{DynamicGameObjectProcessor}
        \begin{itemize}
            \item \textit{DynamicGameObjectEventProcessor}
        \end{itemize}
        \item \textbf{IPropertyProcessor}
    \end{itemize}
\end{frame}


\begin{frame}[fragile]{Rendering: Effekt und Animation}
    \textbf{Partikeleffekt}
    \begin{itemize}
        \item \textbf{Partikel}
        \begin{itemize}
            \item Position
            \item Color
            \item ...\par
        \end{itemize}
        \item \textbf{Partikel Strategy}
        \begin{itemize}
            \item FirePartikelStrategy
            \item JetPartikelStrategy
            \item ...
        \end{itemize}
        \item \textbf{EffectMesh}
        \begin{itemize}
            \item Darstellung im SceneGraph
        \end{itemize}
    \end{itemize}
\end{frame}


\begin{frame}[fragile]{Rendering: Effekt und Animation}
    \textbf{Animationen}
    \begin{itemize}
        \item Steuerung animierter Objekte:
        \begin{itemize}
            \item com.jme3.animation.Animation
            \item com.jme3.animation.AnimChannel
            \item com.jme3.animation.AnimControl
        \end{itemize}
    \end{itemize}
\end{frame}


\begin{frame}[fragile]{Rendering: Pipeline}
    \setlength{\TPHorizModule}{\textwidth}
    \setlength{\TPVertModule}{\textwidth}
    \setbeamercolor{block title}{use=structure,fg=black,bg=white!75!gray}
    \setbeamercolor{block body}{use=structure,fg=black,bg=white!90!gray}
    \begin{itemize}
    \item \begin{textblock}{0.5}(0.05,0.2)
            \begin{block}{Shader}
                \begin{itemize}
                    \item \textbf{VertexShader}
                    \item \textbf{FragmentShader}
                \end{itemize}
            \end{block}
        \end{textblock}
    \item \begin{textblock}{0.5}(0.05,0.5)
            \begin{block}{Postprocessing}
                \begin{itemize}
                    \item \textbf{Motionblur}
                    \item \textcolor{gray}{ShadowFilter}
                \end{itemize}
            \end{block}
        \end{textblock}
    \end{itemize}
\end{frame}


\begin{frame}[fragile]{Rendering: Pipeline}
    \setlength{\TPHorizModule}{\textwidth}
    \setlength{\TPVertModule}{\textwidth}
    \setbeamercolor{block title}{use=structure,fg=black,bg=white!75!gray}
    \setbeamercolor{block body}{use=structure,fg=black,bg=white!90!gray}

    \begin{textblock}{0.35}(0.6,0.4)
        \begin{block}{Materialdefinitionen}
        \mbox{}
        \end{block}
    \end{textblock}
    \begin{itemize}
    \item \begin{textblock}{0.5}(0.05,0.2)
            \begin{block}{Shader}
                \begin{itemize}
                    \item \textbf{VertexShader}
                    \item \textbf{FragmentShader}
                \end{itemize}
            \end{block}
        \end{textblock}
    \item \begin{textblock}{0.5}(0.05,0.5)
            \begin{block}{Postprocessing}
                \begin{itemize}
                    \item \textbf{Motionblur}
                    \item \textcolor{gray}{ShadowFilter}
                \end{itemize}
            \end{block}
        \end{textblock}
    \end{itemize}
\end{frame}


\begin{frame}

    \centering
        {\scshape\LARGE Valaris \par}
    
    \vspace{1cm}

    \includegraphics[height=2.5cm]{./Bilder/ValarisLogo.png}

    \vspace{1cm}
    
    \tiny{\textbf{Quellen:} Alle Bilder aus Prototypen.}

\end{frame}

\end{document}



